\chapter{Profit Maximisation and Supply}\label{ch:b1:profit-maximisation}

The retail price of a pound of white flour in the United States averaged
33 cents in 2006 and 51 cents in 2008, after a rise in world grain prices.
% figdata: flour
American farmers had already moved: the area harvested of maize rose from
28.6 million hectares in 2006 to 35.0 million in 2007, while the area of
soya beans fell from 30.2 to 26.0 million, as land was switched from one
crop to the other. % figdata: crop_area_points
Behind each hectare is a decision about what to grow and with how much
fertiliser, taken by a firm that cannot change the price it gets. This
chapter derives that decision: what a price-taking firm produces at given
prices, how it responds when they change, and whether the model's firm
resembles real ones.

\begin{openingchart}{Area harvested of maize and soya beans, Brazil and the
United States, 1961--2024, millions of hectares.}{FAOSTAT, Crops and
livestock products: area harvested}
\begin{tikzpicture}
\begin{axis}[omtime, xlabel={Year}, ylabel={Million hectares}, xmin=1961, xmax=2024, ymin=0, ymax=50,
  width=0.72\linewidth, height=0.44\linewidth, legend pos=north west]
  \addplot[demand, mark=none] table[x=year, y=maize_usa, col sep=comma] {figdata/bachelor-1/11-profit-maximisation/out/crop_area.csv};
  \addplot[demand, dashed, mark=none] table[x=year, y=soya_usa, col sep=comma] {figdata/bachelor-1/11-profit-maximisation/out/crop_area.csv};
  \addplot[supply, mark=none] table[x=year, y=maize_bra, col sep=comma] {figdata/bachelor-1/11-profit-maximisation/out/crop_area.csv};
  \addplot[supply, dashed, mark=none] table[x=year, y=soya_bra, col sep=comma] {figdata/bachelor-1/11-profit-maximisation/out/crop_area.csv};
  \legend{Maize: United States, Soya: United States, Maize: Brazil, Soya: Brazil}
\end{axis}
\end{tikzpicture}
\end{openingchart}

The chart has fast and slow movements. Year to year, American farmers
switch hectares between maize and soya beans as their relative prices
move. Over decades, Brazil's soya area grew from 1.3 million hectares in
1970 to 45.9 million in 2024, as new land was cleared and new varieties
adapted soya to the tropics; it now exceeds the American
area. % figdata: crop_area_points
A firm's supply has a short run, in which land and machines are fixed, and
a long run, in which they are not.

\begin{recall}
A price-taking firm produces where price equals marginal cost, and shuts
down if the price is below its average variable cost: this was met at
school. The \omterm{def:b1:cost-minimisation:cost-function}{cost function}, the short and the long run and Shephard's lemma
are in \cref{ch:b1:cost-minimisation}; the envelope theorem is
\cref{thm:b1:comparative-statics:envelope}.
\end{recall}

\section{The price-taking firm}\label{sec:b1:profit-maximisation:firm}

A firm sells its output at the price $p$ and hires labour and capital at $w$
and $r$, all given. It chooses its inputs to
\[
\max_{K, L}\ pf(K, L) - wL - rK.
\]
At an interior maximum,
\[
pf_L(K, L) = w,\qquad pf_K(K, L) = r:
\]
each input is hired until the value of its \omterm{def:b1:technology:marginal-product}{marginal product}, $p$ times
the \omterm{def:b1:technology:marginal-product}{marginal product}, equals its price. One more hour of work adds
$f_L$ units worth $pf_L$ and costs $w$. The second-order condition requires
$f$ to be concave near the optimum: diminishing \omterm{def:b1:technology:marginal-product}{marginal products}, and
more.

With constant returns there is no interior maximum. If a plan makes a
profit $\pi > 0$, doubling it makes $2\pi$, so the firm would grow without
limit; if every plan loses money, it makes nothing; and if the best plans
break even, any size will do. Price-taking firms with constant returns
have no determinate size --- a point the industry chapter
(\cref{ch:b1:competitive-industry}) turns into a theory of the number of
firms.

\section{Supply and factor demands}\label{sec:b1:profit-maximisation:supply}

The firm's problem splits in two: for each output, the cheapest way of
making it (\cref{ch:b1:cost-minimisation}), and then the best output,
\[
\max_q\ pq - C(q).
\]
Its first-order condition is $p = C'(q)$: price equals marginal cost. The
second-order condition is $C''(q) \ge 0$: marginal cost must be rising
there.

\begin{definition}[Supply function]\label{def:b1:profit-maximisation:supply-function}
\index{supply function (of a firm)}\index{supply curve!of a firm}
The \emph{supply function} $q^*(p, w, r)$ of a price-taking firm gives its
profit-maximising output at the prices $(p, w, r)$. Its graph against $p$
is the firm's \emph{supply curve}.
\end{definition}

\begin{definition}[Shut-down price]\label{def:b1:profit-maximisation:shut-down}
\index{shut-down price}
In the short run a firm can shut down, losing its fixed cost but paying no
variable cost. It produces only if revenue covers variable cost,
$pq \ge C_v(q)$, that is $p \ge C_v(q)/q$. The \emph{shut-down price} is the
minimum of average variable cost; below it the firm supplies nothing.
\end{definition}

The short-run supply curve is therefore the rising part of marginal cost
above the minimum of average variable cost. In the long run there is no
fixed cost and the firm leaves the industry if the price is below the
minimum of average cost.

\begin{omfigure}
\begin{tikzpicture}
\begin{axis}[omrecord, grid=none, xlabel={Output $q$}, ylabel={Price and cost},
  xmin=0, xmax=25, ymin=0, ymax=14, width=0.6\linewidth, height=0.44\linewidth, samples=100,
  legend pos=south east]
  \fill[omDef!14] (axis cs:0,6.053) rectangle (axis cs:17.971,9);
  \addplot[shifted, domain=4:25] {50/x + 4 - 0.4*x + 0.02*x^2};
  \addplot[guide, domain=0:25] {4 - 0.4*x + 0.02*x^2};
  \addplot[supply, thin, domain=0:10] {4 - 0.8*x + 0.06*x^2};
  \addplot[supply, line width=2.2pt, domain=10:22.2] {4 - 0.8*x + 0.06*x^2};
  \addplot[demand, thin, forget plot] coordinates {(0,9) (25,9)};
  \addplot[supply, line width=2.2pt, forget plot] coordinates {(0,0) (0,2)};
  \addplot[only marks, mark=*, black, forget plot] coordinates {(10,2) (17.971,9)};
  \node[font=\footnotesize, anchor=south west] at (axis cs:0.3,9) {$p = 9$};
  \node[font=\footnotesize] at (axis cs:4.5,7.5) {profit};
  \legend{$AC$, $AVC$, $MC$, supply}
\end{axis}
\end{tikzpicture}

{\small $C(q) = 50 + 4q - 0.4q^2 + 0.02q^3$. The supply curve (thick red) is marginal
cost above the minimum of average variable cost, $(10, 2)$; below the price 2
the firm supplies nothing. At $p = 9$ it makes 17.97 units at an average cost
of 6.05: the shaded profit is $17.97 \times (9 - 6.05) \approx 53.0$ (52.96 unrounded).}
\end{omfigure}

The firm also chooses its inputs; their quantities as functions of all
prices are its demands for them.

\begin{definition}[Unconditional factor demand]\label{def:b1:profit-maximisation:factor-demand}
\index{factor demand!unconditional}\index{unconditional factor demand}
The \emph{unconditional factor demands} $L(p, w, r)$ and $K(p, w, r)$ are the
inputs of the profit-maximising plan. They equal the conditional demands
of \cref{ch:b1:cost-minimisation} at the profit-maximising output:
$L(p, w, r) = L^*(w, r, q^*(p, w, r))$.
\end{definition}

A rise in the wage acts through two channels: at given output the firm
substitutes capital for labour (the conditional demand falls), and since
marginal cost rises the firm makes less, so it needs less of every normal
input. Both push labour demand down.

\section{The profit function}\label{sec:b1:profit-maximisation:profit-function}

\begin{definition}[Profit function]\label{def:b1:profit-maximisation:profit-function}
\index{profit function}
The \emph{profit function} $\pi(p, w, r)$ is the largest profit a
price-taking firm can make at the prices $(p, w, r)$:
$\pi(p, w, r) = \max_{K, L}\ pf(K, L) - wL - rK$.
\end{definition}

Like the \omterm{def:b1:cost-minimisation:cost-function}{cost function}, it collects the firm's whole behaviour in one
function of prices. It is nondecreasing in $p$, nonincreasing in $w$ and $r$,
homogeneous of degree one in all prices together, and convex.

\begin{proposition}[Convexity of the profit function]\label{prop:b1:profit-maximisation:convex}
$\pi(p, w, r)$ is convex in $(p, w, r)$.
\end{proposition}

\begin{proof}
Write a plan as $y = (q, -L, -K)$ and a price vector as $P = (p, w, r)$, so the
profit of $y$ is $P\cdot y$. Let $P = tP^1 + (1 - t)P^2$ and let $y$ be the best plan at
$P$. Then $\pi(P) = tP^1\cdot y + (1 - t)P^2\cdot y \le t\pi(P^1) + (1 - t)\pi(P^2)$, since $y$ is
feasible but not necessarily best at $P^1$ and $P^2$.
\end{proof}

Convexity has the same economic content as the concavity of the \omterm{def:b1:cost-minimisation:cost-function}{cost
function}: a firm that kept its plan when the price rose would see profit
rise along a straight line; by adjusting it does better, so the \omterm{def:b1:profit-maximisation:profit-function}{profit
function} lies above its tangent.

\begin{theorem}[Hotelling's lemma]\label{thm:b1:profit-maximisation:hotelling}
\index{Hotelling's lemma}
Where the supply and factor demands are differentiable,
\[
\frac{\partial\pi}{\partial p} = q^*(p, w, r),\qquad \frac{\partial\pi}{\partial w} = -L(p, w, r),\qquad
\frac{\partial\pi}{\partial r} = -K(p, w, r).
\]
\end{theorem}

\begin{proof}
By the envelope theorem the derivative of the maximum value with respect to
a parameter is the partial derivative of the objective,
$pf(K, L) - wL - rK$, at the optimum: $f(K, L) = q^*$ for $p$, $-L$ for $w$, $-K$
for $r$.
\end{proof}

\begin{proposition}[Law of supply]\label{prop:b1:profit-maximisation:law-of-supply}
\index{law of supply}
A price-taking, profit-maximising firm's supply never falls when its
output price rises, and its demand for an input never rises when the
input's price rises: $\partial q^*/\partial p \ge 0$, $\partial L/\partial w \le 0$, $\partial K/\partial r \le 0$.
\end{proposition}

\begin{proof}
By Hotelling's lemma, $\partial q^*/\partial p = \partial^2\pi/\partial p^2$ and
$\partial L/\partial w = -\partial^2\pi/\partial w^2$; the second derivatives of a convex function
along each axis are nonnegative.
\end{proof}

The law of supply needs no assumption about the shape of the \omterm{def:b1:technology:production-function}{production
function} beyond the existence of a best plan: it follows from profit
maximisation alone. The same argument with two price vectors and two
plans, without derivatives, gives $(p^1 - p^0)(q^1 - q^0) \ge 0$
(\cref{exo:b1:profit-maximisation:convex}).

\begin{omfigure}
\begin{tikzpicture}
\begin{axis}[omrecord, grid=none, xlabel={Output price $p$}, ylabel={Profit},
  xmin=0, xmax=4, ymin=-2, ymax=16, width=0.5\linewidth, height=0.4\linewidth, samples=60,
  legend pos=north west]
  \addplot[demand, domain=0:4] {x^2};
  \addplot[shifted, domain=0.4:4] {4 + 4*(x - 2)};
  \addplot[only marks, mark=*, black, forget plot] coordinates {(2,4)};
  \legend{$\pi(p) = p^2$, plan of $p = 2$ kept}
\end{axis}
\end{tikzpicture}

{\small A firm with $C = q^2/4$ supplies $q = 2p$ and earns $\pi = p^2$. If it kept the
plan chosen at $p = 2$ ($q = 4$), its profit would move along the dashed line;
adjusting, it does better at every other price, so the \omterm{def:b1:profit-maximisation:profit-function}{profit function} is
convex and its slope at $p$ is the output, $\pi'(p) = 2p$.}
\end{omfigure}

\begin{omchart}{Cereal yield, five regions, 1961--2024, tonnes per
hectare.}{FAOSTAT, Crops and livestock products: yield of cereals, primary}
\begin{tikzpicture}
\begin{axis}[omtime, xlabel={Year}, ylabel={Tonnes per hectare}, xmin=1961, xmax=2024, ymin=0, ymax=7,
  width=0.72\linewidth, height=0.42\linewidth, legend pos=north west]
  \addplot+[no marks, thick] table[x=year, y=americas, col sep=comma] {figdata/bachelor-1/11-profit-maximisation/out/cereal_yield.csv};
  \addplot+[no marks, thick] table[x=year, y=asia, col sep=comma] {figdata/bachelor-1/11-profit-maximisation/out/cereal_yield.csv};
  \addplot+[no marks, thick] table[x=year, y=europe, col sep=comma] {figdata/bachelor-1/11-profit-maximisation/out/cereal_yield.csv};
  \addplot+[no marks, thick] table[x=year, y=world, col sep=comma] {figdata/bachelor-1/11-profit-maximisation/out/cereal_yield.csv};
  \addplot+[no marks, thick] table[x=year, y=africa, col sep=comma] {figdata/bachelor-1/11-profit-maximisation/out/cereal_yield.csv};
  \legend{Americas, Asia, Europe, World, Africa}
\end{axis}
\end{tikzpicture}
\end{omchart}

Supply grows on the \omterm{def:b1:labour-supply-and-saving:margins}{intensive margin} too. The world's cereal yield rose
from 1.35 to 4.21 tonnes per hectare between 1961 and 2024, $1.8\%$ a year;
Africa's from 0.81 to 1.68, $1.2\%$ a year. % figdata: yield_growth
Higher yields came from better seeds, more fertiliser and irrigation: a
change in the \omterm{def:b1:technology:production-function}{production function} (\omterm{def:b1:technology:technical-progress}{technical progress}) as much as a
response to prices along it.

\begin{measure}[Area harvested and yield]
FAOSTAT's area harvested counts the hectares from which a crop was
gathered in the year; its production is the quantity gathered; yield is
production divided by area harvested. A yield can rise because each
hectare produces more, because fewer harvests fail, or because poor land
leaves the crop; it falls when new, poorer land is brought in. A rising
area and a rising yield are the extensive and the \omterm{def:b1:labour-supply-and-saving:margins}{intensive margin} of a
farm sector's supply.
\end{measure}

\section{Do firms maximise profit?}\label{sec:b1:profit-maximisation:evidence}

The model does not say that managers solve the equations of this chapter.
It says that firms behave \emph{as if} they did, and it is judged by its
predictions: supply that rises with expected prices, factor demands that
fall with their own price. Two arguments support it. Firms that leave
much profit unearned are less likely to survive competition or to attract
capital; and most decisions are repeated, so a firm can find good rules by
trial. Neither argument is a proof. Survival is slow when competition is
weak, and the record below shows firms that had left large gains unused.

\begin{record}[Management as a missing input]
In a field experiment, free management consulting was given to randomly
chosen plants of large textile firms in India. The plants adopted
practices for quality control, inventory and maintenance, and their
productivity rose by $17\%$ in the first year; within three years the
treated firms had opened more plants. The authors conclude that the main
reason the practices had not been adopted before was that the owners did
not know about them or doubted their value (Bloom, Eifert, Mahajan, McKenzie
and Roberts 2013). % ledger: b1:bloom2013
Profitable improvements had been left unused, and competition had not
forced the badly managed firms out, because limits on managerial time
held back the movement of business between firms. % ledger: b1:bloom2013
\end{record}

\begin{model}{The profit-maximising firm}
\assumes A price-taking firm with a convex technology chooses its output
and inputs to maximise profit.
\predicts Supply that rises with the output price, with a larger response
in the long run than in the short run; factor demands that fall with their
own price; the symmetry $\partial L/\partial r = \partial K/\partial w$.
\tested Planted area responds to expected prices: one estimate of the
supply elasticity of the world's calories from maize, wheat, rice and soya
beans lies between $0.08$ and $0.13$ across specifications (Roberts and
Schlenker 2010), small but not zero. % ledger: b1:roberts2010
The switch from soya to maize in 2007 is the same response in one
country. % figdata: crop_area_points
\breaks Firms that do not know their costs or the best practices
(Bloom et al.\ 2013); % ledger: b1:bloom2013
firms large enough to move their price (\cref{ch:b1:monopoly-pricing}).
\end{model}

\section*{Exercises}\label{sec:b1:profit-maximisation:exercises}

\begin{exercise}[$\star$]\label{exo:b1:profit-maximisation:supply-q2}
A price-taking firm has $C(q) = q^2/4$. Derive its supply function and its
\omterm{def:b1:profit-maximisation:profit-function}{profit function}.
\end{exercise}

\begin{exercise}[$\star$]\label{exo:b1:profit-maximisation:shutdown}
Find the \omterm{def:b1:profit-maximisation:shut-down}{shut-down price} of a firm with
$C(q) = 50 + 10q - q^2 + q^3/3$.
\end{exercise}

\begin{exercise}[$\star$]\label{exo:b1:profit-maximisation:read}
From the opening chart's table, Brazil harvested soya beans from 11.487
million hectares in 1990 and 45.907 million in 2024. % figdata: crop_area_points
By what factor did the area grow? What part of this growth would a supply
curve with a fixed technology and fixed land explain?
\end{exercise}

\begin{exercise}[$\star$]\label{exo:b1:profit-maximisation:sign}
A price-taking firm's wage rises. What is the sign of the change in its
demand for labour, and which result says so?
\end{exercise}

\begin{exercise}[$\star\star$]\label{exo:b1:profit-maximisation:cd}
For $f = K^{1/4}L^{1/2}$, derive the supply and the factor demands, and evaluate
them at $p = 8$, $w = r = 1$.
\end{exercise}

\begin{exercise}[$\star\star$]\label{exo:b1:profit-maximisation:hotelling-check}
Show that the \omterm{def:b1:profit-maximisation:profit-function}{profit function} of the previous exercise is
$\pi = p^4/(64rw^2)$, and check Hotelling's lemma for $q$ and for $L$ at $p = 8$,
$w = r = 1$.
\end{exercise}

\begin{exercise}[$\star\star$]\label{exo:b1:profit-maximisation:predict}
With a supply elasticity of $0.1$, predict the response of the area planted
to a $30\%$ rise in the expected price.
\end{exercise}

\begin{exercise}[$\star\star$]\label{exo:b1:profit-maximisation:yield}
From the second chart's table, cereal yields were 1.353 and 4.206 tonnes per
hectare for the world in 1961 and 2024, and 0.81 and 1.676 for
Africa. % figdata: yield_growth
Compute the yearly growth rates.
\end{exercise}

\begin{exercise}[$\star\star$]\label{exo:b1:profit-maximisation:crs}
Show that a price-taking firm with constant returns and the \omterm{def:b1:cost-minimisation:cost-function}{cost function}
$C(q) = cq$ either makes nothing, wants to grow without limit, or is
indifferent among all sizes. Which happens at which prices?
\end{exercise}

\begin{exercise}[$\star\star\star$]\label{exo:b1:profit-maximisation:convex}
Without derivatives: let the firm choose $q^0$ at the price $p^0$ and $q^1$ at
$p^1$, with input prices fixed. Show that $(p^1 - p^0)(q^1 - q^0) \ge 0$.
\end{exercise}

\begin{exercise}[$\star\star\star$]\label{exo:b1:profit-maximisation:fixed-factor}
A firm has $q = Ax^a\bar K^b$ with $0 < a < 1$, a variable input $x$ at price $v$ and
a fixed $\bar K$. Derive its short-run supply and show that its elasticity is
$a/(1 - a)$.
\end{exercise}

\begin{exercise}[$\star\star\star$]\label{exo:b1:profit-maximisation:refute}
What evidence would refute profit maximisation as a model of farm supply?
Name one pattern in planted areas, and one in input use.
\end{exercise}

\section{Problem: A Kestrian Farm in a Price Boom}\label{sec:b1:profit-maximisation:problem}

\begin{problem}[{Weekend problem --- short-run and long-run supply}]\label{pb:b1:profit-maximisation:farm}
A Kestrian maize farm (invented) produces $q = 40N^{1/4}T^{1/4}$ tonnes from $N$
kilograms of fertiliser, at $v = 2$ kests a kilogram, and $T$ hectares of land,
rented at $\rho = 200$ kests a hectare a year. The farmer's own labour is fixed
and not counted: profit is its return. Maize sells at $p = 100$ kests a tonne.

\medskip\noindent\textbf{Part I --- The plan.}
\begin{enumerate}
  \item What are the \omterm{def:b1:technology:returns-to-scale}{returns to scale}? Why can the farm have a determinate size?
  \item Write the first-order conditions.
  \item Show that the farm spends a quarter of its revenue on each input, and find $q$, $N$ and $T$.
  \item Compute the profit, and the yield per hectare.
  \item Show that the \omterm{def:b1:profit-maximisation:profit-function}{profit function} is $\pi = 1600p^2/(8\sqrt{v\rho})$, and evaluate it.
  \item Check Hotelling's lemma for $q$.
\end{enumerate}

\medskip\noindent\textbf{Part II --- Supply and demands.}
\begin{enumerate}[resume]
  \item Derive the long-run supply function $q(p, v, \rho)$.
  \item What is its price elasticity?
  \item Derive the demand for fertiliser $N(p, v, \rho)$.
  \item What is the elasticity of fertiliser demand with respect to its own price?
  \item Check that $\partial N/\partial\rho = \partial T/\partial v$.
  \item Why does the farm not double its land when it is profitable?
\end{enumerate}

\medskip\noindent\textbf{Part III --- A price boom.} The price of maize rises to 150.
In the short run the land is fixed at 250 hectares.
\begin{enumerate}[resume]
  \item Write the short-run \omterm{def:b1:technology:production-function}{production function} in $N$ alone.
  \item Derive the short-run supply and its elasticity.
  \item Compute short-run output and fertiliser at $p = 150$.
  \item Compute long-run output, land and fertiliser at $p = 150$.
  \item Compare short-run and long-run profit at $p = 150$.
  \item Why is the long-run response larger?
\end{enumerate}

\medskip\noindent\textbf{Part IV --- Fertiliser too.} At the same time the price of
fertiliser rises from 2 to 4.5 kests a kilogram.
\begin{enumerate}[resume]
  \item Compute long-run output.
  \item Compute short-run output, with 250 hectares.
  \item Compute fertiliser per hectare in the short run, against 100 before.
  \item Compute land and fertiliser in the long run.
  \item Explain why output falls in the short run but not in the long run.
  \item State the result: the farm's long-run and short-run supply elasticities.
\end{enumerate}
\end{problem}
